% sec-common-event: oct13.tex:616-636 (label sec:norm-event)
\paragraph{Paper excerpt.}
\begin{verbatim}
\subsection{Martingale estimates}\label{sec:norm-event}

We collect on one event the martingale bounds used in Sections~\ref{sec:norms}--\ref{sec:outer}. Fix $p>0$, and let $K$ be a constant with $\Rout(n)\leq(K-1)r_n$ on the event of Proposition~\ref{prop:coarse}, so that $D_n\subset B(0,Kr_n)$ there for large~$n$. The process
\begin{equation*}
\mathcal Q_t\coloneqq\sum_{j<t}2X_j\cdot\bigl(X_{j+1}-X_j+\drift I_j\xi(X_j)\bigr)
\end{equation*}
is a martingale. Since every step has squared length one, Euler's relation gives
\begin{equation}\label{eq:quadratic}
|X_n|^2=n-2\drift\sum_{x\in A_n}\gauge(x)+\mathcal Q_n\, .
\end{equation}
Let $\widehat{\mathcal Q}$ be~$\mathcal Q$ stopped when the walk leaves~$B(0,Kr_n)$. Its increments are at most~$Cr_n$, so by Azuma's inequality $|\widehat{\mathcal Q}_n|\leq Cr_n\sqrt{n\log n}$ with probability at least $1-Cn^{-p}$, and on the event of Proposition~\ref{prop:coarse} we have $\widehat{\mathcal Q}_n=\mathcal Q_n$, so that
\begin{equation}\label{eq:quadraticerror}
|\mathcal Q_n|\leq Cr_n\sqrt{n\log n}\, .
\end{equation}

To bound the brackets~$\langle\mathcal M^y\rangle_n$ in terms of~$\ell_n(y)$, for $y\in\Z^d$ let $f_y\coloneqq(g_y-\gamma)^2$, where $\gamma\coloneqq0$ if $d\geq3$ and $\gamma\coloneqq(\nf2\pi)\log r_n$ if $d=2$. If $|y|\leq3n$, the increments of~$\mathcal M^{f_y}$ through time~$n$ are at most~$C\bar g$, where $\bar g\coloneqq1$ if $d\geq3$ and $\bar g\coloneqq\log n$ if $d=2$, and Lemma~\ref{lem:freedman} applies to $\mathcal M^{f_y}/(C\bar g)$; we take a union bound over these~$y$.

The event of Section~\ref{sec:norm-event} is the intersection of the events of Proposition~\ref{prop:coarse} and Lemma~\ref{lem:local} with those on which these bounds on~$\widehat{\mathcal Q}_n$ and~$\mathcal M^{f_y}_n$, \eqref{eq:localmart}, and~\eqref{eq:distance-crossing} with $N=n$ hold. The complement of this event has probability at most $Cn^{-p}$, and on the event the arguments of Sections~\ref{sec:norms}--\ref{sec:outer} are deterministic. On the event, $\max_x\ell_n(x)\leq Cr_n$, so~\eqref{eq:approx} gives, for $y\in\R^d$ with $|y|\leq2n$,
\begin{equation}\label{eq:norm-approx}
|\widetilde\ell_n(y)-U_{D_n}(y)|\leq C\begin{cases}\sqrt{r_n}\log n,&d=2,\\\sqrt{r_n\log n},&d\geq3\end{cases}\, .
\end{equation}
\end{verbatim}
\paragraph{Lean statement.}
\begin{verbatim}
theorem CERW.LimitShapes.exists_common_event {d : ℕ} (hd : 2 ≤ d)
    {Ψ : EuclideanSpace ℝ (Fin d) → ℝ} (hΨ : CERW.IsNorm Ψ) {ε : ℝ} (hε : 0 < ε)
    (hell : ∀ i : Fin d, ε * Ψ (CERW.coordVec i) < 1 / (d : ℝ)) (p : ℝ) (hp : 0 < p) :
    ∃ K : CERW.LimitShapes.EventConstants,
      0 < K.cCo ∧ 0 < K.CCo ∧ 0 < K.CLoc ∧ 0 < K.CQuad ∧
      0 < K.CSquare ∧ 0 < K.CMart ∧ 0 < K.CCross ∧
      ∃ C : ℝ, 0 < C ∧ ∃ n₀ : ℕ, 2 ≤ n₀ ∧
        ∀ ξ : Site d → EuclideanSpace ℝ (Fin d),
          (∀ x : Site d, x ≠ 0 → CERW.IsSubgradient Ψ (CERW.toSpace x) (ξ x)) → ξ 0 = 0 →
        ∀ {Ω : Type u} [MeasurableSpace Ω] (μ : Measure Ω) [IsProbabilityMeasure μ]
          (X : ℕ → Ω → Site d), CERW.IsDriftCERW μ ε ξ X → ∀ n : ℕ, n₀ ≤ n →
          μ (CERW.LimitShapes.section5Event Ψ ε ξ K X n)ᶜ ≤
            ENNReal.ofReal (C * (n : ℝ) ^ (-p))
\end{verbatim}
\paragraph{Proof of the registered statement.}
Fix $d,\kappa,\Psi,p$. Obtain the positive coarse constants from node
\texttt{oct13-prop-coarse}, the reduced approximation constant from node
\texttt{oct13-lem-local}, and the two actual-bracket Freedman constants from
\texttt{KernelEvents.exists\_localMartingaleBound} and
\texttt{KernelEvents.exists\_squareMartingaleBound}. Obtain the positive
crossing constant from node \texttt{eq-distance-crossing}. These choices
are made before the subgradient field or law.

For the quadratic constituent, subgradient duality bounds every drift vector by
$\Lambda=\operatorname{normMax}(\Psi)$, uniformly over the subgradient field.
The actual full-prefix stopped process is a martingale, with increments bounded
almost surely by $2R(1+\kappa\Lambda)$. Normalize by this positive bound and use
the proved martingale clamp to obtain a surely bounded, almost surely equal
representative. The existing logarithmic interval estimate at $[0,n]$, with
exponent $p+1$, gives failure at most $4n^{-p}$; transferring through the
almost-sure equality and enlarging the constant proves the direct stopped tail.
Its constant is chosen before the field, law, radius and horizon.

Choose $R=(C_{\rm co}+1)r_n>0$ at the deterministic coarse scale and take
$K.C_{\rm Quad}=C_Q(C_{\rm co}+1)$. This makes the direct tail threshold
exactly the quadratic threshold in the literal event. No unstopped estimate
is assumed to produce this constituent. A deterministic threshold $r_n\geq1$
is retained, and the final threshold is at least two.

Take the sixfold intersection in the literal definition of
\texttt{section5Event}. Failure of
any constituent is contained in the union of the coarse, reduced-local,
quadratic, two kernel-Freedman, or crossing exceptional events. The union
bound gives a coefficient equal to the sum of their positive coefficients.
Enlarge the common time threshold to at least two. This gives positive
constants in every field of the existential witness and the asserted
polynomial measure bound, uniformly in every subgradient selection and
realization. No older seven-constituent event is substituted.
\paragraph{Lean status.}
The registered statement is sealed by applying its proved common-event provider.
